\section{Náhled pro sociální sítě}
Významným zdrojem návštěvnosti stránek mohou v dnešní době být i sociální sítě. Aby byl sdílený odkaz na stránku dostatečně zajímavý je třeba stránky rozšířit o patřičné meta informace. Facebook a většina sociálních sítí ke specifikování náhledu využívá protokol Open Graph představený firmou Facebook v roce 2010 \autocite{opengraph}.

\begin{listing}[htbp]
  \begin{minted}{twig}
{% macro meta(title,url,description,image,imageHeight,imageWidth) %}
    <meta property="og:title" content="{{ title }}">
    <meta property="og:description" content="{{ description }}">
    <meta property="og:url" content="{{ url }}">
    <meta property="og:image" content="{{ image }}">
    <meta property="og:image:width" content="{{ imageWidth }}">
    <meta property="og:image:height" content="{{ imageHeight }}">
    <meta property="og:site_name" content="{{ name }}">

    <meta property="twitter:card" content="summary">
{% endmacro %}
  \end{minted}
  \caption{\label{opengraph-meta}}
\end{listing}

\begin{figure}[htbp]
        \centering
        \includegraphics[width=\textwidth,height=\textheight,keepaspectratio]{images/fbshare.png}
        \caption{\label{facebook-share}}
\end{figure}